\chapter{Zahlendarstellungen}\label{ch:zahlendarstellungen}

Im vorigen Kapitel haben wir gesehen, dass sich die Informatik mit der Darstellung, Verarbeitung und Übertragung von Informationen befasst. Bevor wir uns der Hardware zuwenden, mit der ein Computer diese Informationen tatsächlich verarbeitet, müssen wir uns zunächst mit einer grundlegenden Frage beschäftigen: Wie stellt man Zahlen überhaupt dar? Diese Frage ist keineswegs so selbstverständlich, wie sie zunächst klingt, und ihre Beantwortung wird uns direkt zum \tib{Binärsystem} führen -- der Sprache, in der Computer intern \enquote{denken}.

\section{Ursprünge der Zahlendarstellung}
Die \textit{natürlichen Zahlen} werden so genannt, da sie auf \enquote{natürliche Weise} beim Zählen verwendet werden.
Auf dem Gebiet der heutigen Demokratischen Republik Kongo wurde im 20. Jahrhundert ein Knochen, der sogenannte \textit{Ishango-Knochen}, gefunden.
Dieser Knochen wird auf die Zeit vor etwa 20'000 Jahren vor unserer Zeitrechnung datiert.
In den Knochen wurden deutlich erkennbar von Menschen einige Kerben eingeritzt.
Die Bedeutung dieser Kerben ist unklar, es wird jedoch angenommen, dass sie eine bestimmte Anzahl festhalten sollten.

\begin{figure}[H]
    \centering
    \includegraphics[width=.30\textwidth,height=55mm,keepaspectratio]{Grundlagen_Info/13_ZahlendarstellungenUndKodierungen/Figures/ishango-bone}
    \caption{Der Ishango-Knochen (ca. 20'000 Jahre alt) mit eingeritzten Kerben als historisches Beispiel der Zahlendarstellung.}
    \label{fig:ishango_bone_skript}
\end{figure}

Stellen wir uns vor, dass die Kerben die Anzahl der gefangenen Fische darstellen.
Gefangene Fische durch Kerben in einem Knochen zu repräsentieren, verlangt bereits einen recht hohen Grad an
Abstraktion --- schliesslich haben gefangene Fische und Kerben in einem Knochen auf den ersten Blick keinen offensichtlichen Zusammenhang.
Weniger abstrakt wäre es, eine Darstellung zu wählen, welche direkt an das zu zählende Objekt gebunden ist.
In diesem Fall würde man also keine Kerben ritzen, sondern die entsprechende Anzahl Fische zeichnen / skizzieren.
Eine Anzahl von drei Fischen könnte also durch drei Fisch-Symbole $\twemoji{tropical fish}\twemoji{tropical fish}\twemoji{tropical fish}$
festgehalten werden.

Anstelle von Kerben in einem Knochen oder Fisch-Symbole als Markierung könnte eine bestimmte Anzahl Fische auch durch
die entsprechende Anzahl von anderen Markierungen wie Striche ($|$) oder Kieselsteine (\textbullet) repräsentiert werden:
\begin{align*}
	\text{drei Kerben} \quad\leftrightarrow\quad \twemoji{tropical fish}\twemoji{tropical fish}\twemoji{tropical fish} \quad\leftrightarrow\quad |||
	\quad\leftrightarrow\quad \text{\textbullet\textbullet\textbullet}
\end{align*}
Die Darstellung durch skizzierte Fische ist am wenigsten abstrakt, da die dabei verwendete Markierung den zu zählenden Objekten (in diesem Fall also den Fischen) direkt nachempfunden ist.
Zahlendarstellungen, welche nur ein Symbol / eine Markierung (Kerbe, \twemoji{tropical fish}, $|$, \textbullet) verwenden, werden
\tib{unäre Zahlendarstellungen}\index{Zahlendarstellung!unäre} genannt.

\begin{myremark}\leavevmode
	\begin{enumerate}[(a)]
		\item Der Vorteil von unären Darstellungen (gegenüber \enquote{komplexeren} Darstellungen) ist ihre einfache Verständlichkeit und
		      offensichtliche Interpretation.
		\item Unäre Darstellungen sind nur für kleine natürliche Zahlen übersichtlich, da ihre Darstellungslängen linear mit der Grösse der zu
		      darstellenden Zahl wachsen (doppelt so grosse Zahl $\rightarrow$ doppelt so lange Darstellung).
		\item Unäre Darstellungen eignen sich nicht besonders gut zum Rechnen.
	\end{enumerate}
\end{myremark}

\section{Basisgrössen}\label{section:Basisgroessen}
Anstelle von nur einem einzigen Symbol, wie bei unären Darstellungen, haben bereits Kulturen der Antike unterschiedliche
\tib{Basisgrössen}\index{Basisgrössen} verwendet. \enquote{Basisgrössen} bedeuten in diesem Kontext für Symbole, die gewisse Mengen, bzw. Zahlen repräsentieren, und dies unabhängig vom Objekt, das gezählt wird. Die Römer verwendeten die sieben verschiedenen Basisgrössen:
\begin{center}
	$I$ (Eins), $V$ (Fünf), $X$ (Zehn), $L$ (Fünfzig), $C$ (Hundert), $D$ (Fünfhundert) und $M$ (Tausend).
\end{center}
\begin{myexample}\label{myexample:roman}
Die Zahl \textit{fünfhundertvierundfünfzig} könnte in den römischen Basisgrössen zum Beispiel durch
\begin{align*}
	CCCCCLIIII
\end{align*}
oder durch
\begin{align*}
	DLIIII
\end{align*}
dargestellt werden.
\end{myexample}
Wir werden uns jede Basisgrösse als einen Typ von Münze von einem bestimmten Wert vorstellen. Die darzustellende Zahl
denken wir uns dann als einen bestimmten Geldbetrag, den wir mithilfe dieser gegebenen Basisgrössen (Münztypen) auszahlen wollen. In
dieser Vorstellung hatten die Römer also ein Währungssystem von sieben verschiedene Münztypen. Der Geldbetrag \textit{einundzwanzig} könnte
dann beispielsweise durch $XXI$ (zwei Zehnermünzen und eine Einermünze) in antiker römischer Währung ausbezahlt werden. Die alten Ägypter hingegen verwendeten damals bereits Symbole (Zeichnungen) für die Basisgrössen $1, 10, 100, 1000, 10000,
	100000$ und $1000000$.

Ein Geldbetrag kann sowohl im römischen als auch ägyptischen System typischerweise auf mehrere unterschiedliche Arten
ausbezahlt werden (zum Beispiel, indem nur Einer verwendet werden). Sie kennen diese Situation vom Bezahlen mit
Bargeld --- so kann der Betrag von hundert Schweizer Franken auf viele verschiedene Arten mit Schweizer Bargeld bezahlt werden
(z.B. mit zwei Fünfzigernoten, fünf Zwanzigernoten. Natürlich ist
es bei Zahlendarstellungen vorteilhaft, wenn sie \textit{eindeutig} sind. Dies bedeutet, dass es für jede natürliche Zahl eine einzige minimale Münzenmenge gibt, die diese Zahl repräsentiert. Für digitale Systeme ist dies essenziell: Ein Computer benötigt für jeden Wert ein \textbf{eindeutiges Bitmuster} (damit Werte ohne Mehrdeutigkeit verglichen und gespeichert werden können) sowie eine \textbf{minimale Stellenanzahl} (um Speicherplatz zu sparen). Angenommen wir haben einen Betrag $x$ gegeben, welcher durch
unterschiedliche Ansammlungen von Münzen ausbezahlt werden kann. Dann betrachten wir diejenige Ansammlung von Münzen als
die \tib{Zahlendarstellung}\index{Zahlendarstellung} von $x$, welche am wenigsten Münzstücke benötigt.

\begin{myexample}\label{myexample:notUnique}
Der Betrag \textit{sieben} kann im römischen System auf genau zwei verschiedene Arten (abgesehen von der Reihenfolge) ausbezahlt werden:
\begin{itemize}
	\item $IIIIIII$ (benötigt sieben Münzstücke),
	\item $VII$ (benötigt drei Münzstücke).
\end{itemize}
Die Darstellung $VII$ verwendet dabei am wenigsten Münzstücke und wird somit als die \textit{minimale} Zahlendarstellung von
\textit{sieben} verstanden.
\end{myexample}
Nun gibt es aber denkbare Münzsysteme, in denen nicht jeder Betrag auf eindeutige Weise mit möglichst wenigen
Münzstücken ausbezahlt werden kann. Dies werden Sie in \cref{myexercise:1001} untersuchen.

\begin{myexercise}\label{myexercise:1001}
	Erweitern Sie das römische Münzsystem um zwei weitere Münztypen Ihrer Wahl und finden Sie einen Betrag, der in Ihrem
	neuen System mehrere unterschiedliche Darstellungen mit derselben minimalen Anzahl von Münzen erlaubt.
	\begin{myanswer}
		Wir erweitern das römische Münzsystem um die beiden Münztypen \coin{II} und \coin{III}. Dann besitzt der Betrag
		\textit{vier} die beiden Darstellungen \coin{I}\coin{III} oder \coin{II}\coin{II}, welche beide eine (minimale) Anzahl von zwei Münzstücken verwenden.
	\end{myanswer}
\end{myexercise}


\begin{myexercise}\label{myexercise:grosse-roem-zahlen}
	Ein Haupt-Nachteil von Zahlendarstellungen durch Münzsysteme ist, dass Münzsysteme für grosse Zahlen zur gleichen Ineffizienz führen wie unäre Zahlensysteme. Wie viele Zeichen würde man benötigen, um die Zahl \num{1000000000} (1 Milliarde) im römischen Zahlensystem darzustellen?
	\begin{myanswer}
		Die Zahl wäre eine Million Stellen lang: Die Zahl \textit{M} (Tausend) ist die grösste Zahl des römischen Zahlensystems und wir benötigen eine Million mal Tausend, um auf eine Milliarde zu kommen.
	\end{myanswer}
\end{myexercise}
\section{Stellenwertdarstellung}
\label{section:stellenwertdarst}
Das uns vermutlich am besten vertraute Zahlensystem ist das \tib{Dezimalsystem}\index{Stellenwertsystem!Dezimalsystem}
(Zehnersystem). Es verwendet die \textbf{zehn} verschiedenen Ziffern
\begin{align*}
	\lrc{0,1,2,3,4,5,6,7,8,9}
\end{align*}
und die (natürlichen) Potenzen $10^0, 10^1, 10^2, \ldots$ der Basis $10$ als \tib{Stellenwerte}\index{Stellenwert}.
Bereits in der Primarschule haben Sie gelernt, dass beispielsweise das Symbol $\green{4}\orange{6}\blue{0}\red{5}$ definiert ist als die folgende
Summe:
\begin{center}
	\red{fünf} Einer ($\red{5}\cdot 10^0$) $+$ \blue{null} Zehner ($\blue{0}\cdot 10^1$) $+$ \orange{sechs} Hunderter
	($\orange{6}\cdot 10^2$) + \green{vier} Tausender
	($\green{4}\cdot 10^3$).
\end{center}
Eine Ziffer gibt dabei jeweils an, wie viele Male die entsprechende Potenz der Basis $10$ verwendet werden soll. Das
Dezimalsystem hat also (die unendlich vielen verschiedenen) Basisgrössen
\begin{align*}
	10^0=1,\quad 10^1=10,\quad 10^2=100,\quad 10^3=1000,\ldots
\end{align*}

\begin{mydefinition}[Stellenwertdarstellung]\label{mydefinition:Stellenwertdarstellung}
Seien $a_0, a_1, \ldots, a_n$ natürliche Zahlen. Dann ist der Ausdruck
\begin{align*}
	a_na_{n-1}\ldots a_1a_0
\end{align*}
definiert durch die Summe
\begin{align*}
	a_n\cdot 10^n + a_{n-1}\cdot 10^{n-1}+\ldots + a_1\cdot 10^1 + a_0\cdot 10^0.
\end{align*}
\end{mydefinition}

\begin{myexample}
Der Ausdruck $4205$ entspricht der Zahl
\begin{align*}
	\textcolor{red}{4205} := \textcolor{red}{\overbrace{4}^{a_3}}\cdot \underbrace{10^3}_{=1000} + \textcolor{red}{\overbrace{2}^{a_2}}\cdot \underbrace{10^2}_{=100} + \textcolor{red}{\overbrace{0}^{a_1}}\cdot \underbrace{10^1}_{=10} + \textcolor{red}{\overbrace{5}^{a_0}}\cdot \underbrace{10^0}_{=1}.
\end{align*}
\end{myexample}

Anstelle der Potenzen der Basis $10$ kann man auch andere natürliche Zahlen als Basis verwenden.
Das \tib{Binärsystem}\index{Stellenwertsystem!Binärsystem} (Zweiersystem) verwendet zum Beispiel die Basis $2$ und
Ziffern aus $\lrc{0,1}$.
\begin{myexample}
Wir wollen den Ausdruck $1101$ in Basis $\blue{2}$ interpretieren. Um dies zu verdeutlichen, schreibt man die Basis
tiefgestellt dahinter, also $1101_{\blue{2}}$. Der Ausdruck $1101_2$ wird dann aufgefasst als die Zahl
\begin{align*}
	1\cdot 2^3 + 1\cdot 2^2 + 0\cdot 2^1 + 1\cdot 2^0 = 1\cdot 8 + 1\cdot 4 + 0\cdot 2 + 1\cdot 1 = 13.
\end{align*}
Die binäre Zahl $1101_2$ entspricht also der Zahl dreizehn.
\end{myexample}

Das Binärsystem findet insbesondere in der Informatik Anwendung, da es sich gut für die Darstellung von Informationen in elektronischen Schaltungen eignet. In einem elektronischen Schaltkreis können beispielsweise zwei verschiedene Spannungspegel verwendet werden, um die beiden Ziffern 0 und 1 zu repräsentieren. Dadurch können Informationen in Form von Binärzahlen gespeichert und verarbeitet werden. Wie ein solcher Schaltkreis konkret aussieht, und wie er mit den beiden Zuständen \enquote{0} und \enquote{1} rechnet, sehen wir in \cref{ch:logische_gatter}.

\begin{myattention}
	Falls wir nicht explizit eine Basis tiefgestellt hinter einen Zahlenausdruck schreiben, dann ist dieser Ausdruck immer in Basis $10$ (im
	Dezimalsystem) aufzufassen. Somit meinen wir also mit $11$ die Zahl elf und nicht etwa die Zahl $11_2 = 1\cdot 2^1 +
		1\cdot 2^0$.
\end{myattention}

\begin{myexercise}
	Schreiben Sie die Binärzahl $11011_2$ als Dezimalzahl.
	\begin{myanswer}
		\begin{align*}
			11011_2 = 1\cdot 2^4 + 1\cdot 2^3 + 0\cdot 2^2 + 1\cdot 2^1 + 1\cdot 2^0 = 16 + 8 + 2 + 1 = 27_{10}
		\end{align*}
	\end{myanswer}
\end{myexercise}

\begin{myexercise}
	Schreiben Sie die Dezimalzahl $14$ als Binärzahl.
	\begin{myanswer}
		\begin{align*}
			14_{10} = 8 + 4 + 2 = 1\cdot 2^3 + 1\cdot 2^2 + 1\cdot 2^1 + 0\cdot 2^0 = 1110_2
		\end{align*}
	\end{myanswer}
\end{myexercise}

Warum schränken wir uns beim Schreiben von Binärzahlen eigentlich auf den Gebrauch von lediglich den beiden Ziffern aus
$\lrc{0,1}$ ein? Ein Ausdruck der Form
\begin{align*}
	3\cdot 2^2 + 5\cdot 2^1 + 7\cdot 2^0
\end{align*}
wäre doch ebenfalls sinnvoll. Diese Frage wollen wir in den folgenden Aufgaben klären. Erinnern Sie sich zurück an \cref{section:Basisgroessen}.
Dort haben wir versucht, Geldbeträge mit möglichst wenigen Münzstücken auszuzahlen.

\begin{myexercise}
	Der Ausdruck $3\cdot 2^2 + 5\cdot 2^1 + 7\cdot 2^0$ verwendet $3+5+7 = 15$ Münzstücke und stellt die Zahl $29$ dar.
	Stellen Sie die Zahl $29$ im Binärsystem dar, verwenden Sie diesmal aber lediglich Ziffern aus $\lrc{0,1}$. Wie viele
	Münzstücke verwendet diese Darstellung?
	\begin{myanswer}
		\begin{align*}
			29_{10} = 1\cdot 2^4 + 1\cdot 2^3 + 1\cdot 2^2 + 0\cdot 2^1 + 1\cdot 2^0 = 11101_2
		\end{align*}
		Diese Darstellung verwendet lediglich 4 Münzstücke.
	\end{myanswer}
\end{myexercise}

\begin{myexample}\label{myexample:ersetzen}
Wir wollen dreizehn Objekte im Dezimalsystem mit so wenigen Münzstücken wie möglich darstellen. Wir könnten dies mit
dreizehn Einern tun. Doch wir sehen sofort, dass sich zehn Einer stets durch nur einen einzigen Zehner (und null
Einer) ersetzen lassen. Dadurch haben wir insgesamt neun Münzen gespart:
\begin{itemize}
	\item verfügbare Münztypen: \coin{1}, \coin{10}, $\ldots$
	\item Münzen vorher:
	      \coin{1}\coin{1}\coin{1}\coin{1}\coin{1}
	      \coin{1}\coin{1}\coin{1}\coin{1}\coin{1}
	      \coin{1}\coin{1}\coin{1}
	\item Münzen nachher: \coin{10}\coin{1}\coin{1}\coin{1}
\end{itemize}
\end{myexample}

\begin{mydefinition}[$\blue{b}$-adische Zahlendarstellung für natürliche Zahlen]
\label{mydefinition:badisch}
Gegeben ist eine natürliche Zahl $\blue{b}\geq 2$ (\tib{Basis}\index{Basis}). Dann definieren wir
\begin{align*}
	(\red{a}_{n}\red{a}_{n-1}\ldots \red{a}_{1}\red{a}_{0})_{\blue{b}} := \red{a}_n\cdot \blue{b}^n + \red{a}_{n-1}\cdot \blue{b}^{n-1} + \ldots +\red{a}_1\cdot \blue{b}^1 + a_0\cdot
	\blue{b}^0,
\end{align*}
wobei die Zahlen $a_{0}, a_{1},\ldots, a_{n-1}, a_{n}\in\lrc{0,1,\ldots,b-1}$ \tib{Ziffern}\index{Ziffern} genannt werden. Wird
nicht explizit eine andere Basis genannt, so ist stets die Basis $b=10$ gemeint.

Diese Zahlendarstellung wird $b$-\tib{adische Zahlendarstellung} genannt.
\end{mydefinition}

\begin{myexample}
Für die Wahl $b:=3$ erhalten wir die $3$-adische Zahlendarstellung (Dreiersystem). Die Zahl $211_3$ entspricht dann der Zahl
\begin{align*}
	211_3 = 2\cdot 3^2 + 1\cdot 3^1 + 1\cdot 3^0 = 2\cdot 9 + 1\cdot 3 + 1\cdot 1 = 22.
\end{align*}
\end{myexample}

\begin{mychallenge}
	Begründen Sie, warum es nicht besonders sinnvoll ist, die natürliche Zahl $1$ als Basis zu wählen, selbst wenn wir in \cref{mydefinition:badisch} die Einschränkung $a_0, a_1, \ldots, a_{n-1}, a_{n}\in\lrc{0,1,\ldots,b-1}$ an die Ziffern ignorieren?
	\begin{myanswer}
		Mit der Wahl der Basis $b := 1$ erhalten wir keine unterschiedlichen Stellenwerte, da $b^k = 1^k = 1$ für alle natürlichen Zahlen $k$ gilt.
	\end{myanswer}
\end{mychallenge}

\begin{myremark}[Sinn der Zifferneinschränkung $a_k < b$]\label{myremark:ziffern-einschraenkung}
	Warum ist es sinnvoll in \cref{mydefinition:badisch} zu fordern, dass eine $b$-adische Darstellung keine Ziffern grösser als $b-1$ verwenden darf?
	Unsere Argumentation ist ähnlich zu der in \cref{myexample:ersetzen}. Angenommen eine Darstellung
	\begin{align*}
		\red{a}_n\cdot \blue{b}^n + \red{a}_{n-1}\cdot \blue{b}^{n-1} + \ldots +\red{a}_1\cdot \blue{b}^1 + a_0\cdot
		\blue{b}^0
	\end{align*}
	würde eine Ziffer $a_k \geq b$ verwenden. Anders gesagt, würden dann $b$ oder mehr Münzstücke von dem Typ $b^k$ verwendet werden.
	Dann könnten wir aber immer eine Gruppe von $b$ vielen Münzen vom Typ $b^k$ durch jeweils eine einzige Münze vom Typ $b^{k+1}$ ersetzen und würden für jede solche Gruppe $b-1$ Münzstücke einsparen.

	Beispielsweise ist es in der Basis $b = 10$ nicht sinnvoll 23 Hunderter ($23\cdot 10^2$) zu verwenden, da wir stattdessen zwei Tausender und nur drei Hunderter ($2\cdot 10^3 + 3\cdot 10^2$) verwenden könnten.
	Dabei würden wir insgesamt $18$ Münzstücke einsparen.

	Für die Informatik bringt die Bedingung $a_k \in \lrc{0, 1, \ldots, b-1}$ zwei entscheidende Vorteile:
	\begin{enumerate}[(1)]
		\item \textbf{Minimale Darstellung (Speichereffizienz):} Die Darstellung benötigt so wenige Stellen (Bits) wie möglich.
		\item \textbf{Eindeutigkeit (Kanonische Form):} Jede Zahl besitzt exakt \emph{eine} einzige $b$-adische Repräsentation, sodass Computer zwei Zahlenwerte direkt vergleichen und abspeichern können. Das Bündeln von $b$ Einheiten zur nächsten Stelle entspricht ausserdem physikalisch genau dem \textbf{Übertrag} (\textit{Carry}) in elektronischen Rechenschaltungen (vgl. \cref{sec:halbaddierer}).
	\end{enumerate}
\end{myremark}

\begin{myexercise}
	Ordnen Sie die folgenden Zahlen aufsteigend:
	\begin{align*}
		43_{10}, 2201_3, 11101_2, 33_4, 142_5, 111_2.
	\end{align*}
	\begin{myanswer}
		\begin{align*}
			2201_3  & = 2\cdot 3^3 + 2\cdot 3^2 + 0\cdot 3^1 + 1\cdot 3^0 = 73_{10},              \\
			11101_2 & = 1\cdot 2^4 + 1\cdot 2^3 + 1\cdot 2^2 + 0\cdot 2^1 + 1\cdot 2^0 = 29_{10}, \\
			33_4    & = 3\cdot 4^1 + 3\cdot 4^0 = 15_{10},                                        \\
			142_5   & = 1\cdot 5^2 + 4\cdot 5^1 + 2\cdot 5^0 = 47_{10},                           \\
			111_2   & = 1\cdot 2^2 + 1\cdot 2^1 + 1\cdot 2^0 = 7_{10},                            \\
			        & \Rightarrow 111_2 < 33_4 < 11101_2 < 43_{10} < 142_5 < 2201_3
		\end{align*}
	\end{myanswer}
\end{myexercise}

\begin{mychallenge}
	Die nachfolgende Aussage gilt für jede Basis $b$ eines $b$-adischen Systems. Zur Vereinfachung werden wir aber nur die Wahl der Basis $b := 10$ untersuchen.
	Zeigen Sie, dass die $10$-adische Darstellung eindeutig ist. Es gibt also keine zwei unterschiedlichen $10$-adischen Darstellungen $x$ und $y$, welche dieselbe Zahl (Anzahl Objekte) darstellen.
	\begin{myanswer}
		Seien
		\begin{align*}
			x & := a_n\cdot 10^n + a_{n-1}\cdot 10^{n-1} + \ldots + a_k\cdot 10^k + a_{k-1}\cdot 10^{k-1} + \ldots + a_1\cdot 10^1 + a_0\cdot 10^0 \\
			y & := b_n\cdot 10^n + b_{n-1}\cdot 10^{n-1} + \ldots + b_k\cdot 10^k + b_{k-1}\cdot 10^{k-1} + \ldots + b_1\cdot 10^1 + b_0\cdot 10^0
		\end{align*}
		zwei unterschiedliche $10$-adische Darstellungen derselben Zahl.
		Dann müssen sich $x$ und $y$ an mindestens einer Stelle (in einer Ziffer) unterscheiden.
		Sei $10^k$ die grösste Stelle, an der sich $x$ und $y$ unterscheiden. Wir können annehmen, dass $a_k > b_k$ gilt. Dann bewirkt der Unterschied an dieser Stelle
		eine Differenz von mindestens $10^k$. Doch selbst mit der Wahl $a_{k-1} = \ldots = a_1 = a_0 = 0$ und $b_{k-1} = \ldots = b_1 = b_0 = 9$ kann nur noch ein Unterschied von $10^k - 1$ erwirkt werden.
		Dann ist aber die von $x$ dargestellte Zahl um mindestens 1 grösser als die von $y$ dargestellte Zahl.
	\end{myanswer}
\end{mychallenge}

\subsection{Umrechnungen zwischen Darstellungen}
Wir wollen nun sehen, wie verschiedene Darstellungen ineinander umgerechnet werden können.
\begin{myexample}[Darstellung in Basis $5$]\label{myexample:umrechnung}
Wie lautet die $5$-adische Darstellung der Zahl $84$? Wir beginnen damit alle natürlichen Potenzen von $5$
aufzulisten, die kleiner oder gleich $84$ sind:
\begin{align*}
	5^0 = 1, \quad 5^1 = 5, \quad 5^2 = 25.
\end{align*}
Die Potenz $5^3 = 125 > 84$ ist bereits grösser als die gegebene Zahl. Nun versuchen wir möglichst viele der möglichst
grössten Potenz von $5$ zu verwenden, solange diese noch in $84$ passen. Die grösste mögliche Potenz ist $5^2 = 25$ und
$25$ hat \red{dreimal} in $84$ Platz. Es bleibt also noch $84-3\cdot 5^2 = 9$ übrig. Wir verwenden möglichst viele der
möglichst grössten Potenz von $5$, die in $9$ Platz haben und stellen fest, dass $5^1=5$ \blue{einmal} Platz hat. Dann
bleibt noch $9-5^1 = 4$ übrig. Wir verwenden möglichst viele der möglichst grössten Potenz von $5$, die in $4$ Platz
haben. Wir verwenden also \green{viermal} die Potenz $5^0 = 1$. Nun bleibt nichts mehr übrig und die $5$-adische
Darstellung von $84_{10}$ lautet:
\begin{align*}
	84_{10} = \red{3}\blue{1}\green{4}_5.
\end{align*}
\end{myexample}
Da wir in \cref{myexample:umrechnung} immer mit der grössten Basisgrösse begonnen haben und so viele Münzen wie möglich
davon verwendet haben, nennen wir diese Methode \tib{gierig}\index{gierig / greedy} (englisch: \tib{greedy}).

\subsection{Hexadezimalsystem}

Welche Ziffern sollen wir jedoch wählen, sofern die Basis $b > 10$ gewählt wird? Die dezimale Zahl 10 ist ja keine Ziffer, sondern besteht selber aus zwei Ziffern (genau so, wie es das Dezimalsystem vorschreibt).

Im Grunde genommen sind Ziffern wie \enquote{1}, \enquote{2} usw. nichts weiter als Zeichen, die von Menschen definiert worden sind, und mit welchen eine gewisse Menge gemeint ist. Wir können einfach Buchstaben verwenden, um höhere Ziffernwerte zu bezeichnen. Genau dies wird im sogenannten \tib{Hexadezimalsystem}\index{Hexadezimalsystem} gemacht ($b=16$). Die 16 Ziffern im Hexadezimalsystem sind:
\begin{align*}
	\lrc{0, 1, 2, 3, 4, 5, 6, 7, 8, 9, \text{A}, \text{B}, \text{C}, \text{D}, \text{E}, \text{F}},
\end{align*}
wobei beispielsweise die Ziffer $\text{A} = 10_{10}$, $\text{B} = 11_{10}$ bis $\text{F} = 15_{10}$ bezeichnet.

\begin{myremark}[Notwendigkeit und Vorteil im Computer]
Da $16 = 2^4$ gilt, entspricht \textbf{genau eine Hex-Ziffer} einer Gruppe von \textbf{4 Bits} (einem sogenannten \textit{Nibble}). 

Ein Byte ($= 8$ Bits) lässt sich damit exakt durch \textbf{zwei Hexadezimalziffern} darstellen ($11111111_2 = \text{FF}_{16}$). Das Hexadezimalsystem ist daher eine unentbehrliche Notwendigkeit, um Binärdaten, Speicheradressen und Farbwerte für Menschen kompakt und lesbar zu notieren.
\end{myremark}

\begin{myexercise}
	Wie lautet die Zahl $27_{10}$ in der $16$-adischen Darstellung?
	\begin{myanswer}
		\begin{align*}
			27_{10} = 1\cdot 16^1 + 11\cdot 16^0 = 1\text{B}_{16}
		\end{align*}
	\end{myanswer}
\end{myexercise}

\begin{myexercise}
	Geben Sie folgende Zahlen im hexadezimalen System aus: 44, 23, 90, 124, 306.
	\begin{myanswer}
		$44_{10} = 2\text{C}_{16},\ 23_{10} = 17_{16},\ 90_{10} = 5\text{A}_{16},\ 124_{10} = 7\text{C}_{16},\ 306_{10} = 132_{16}$
	\end{myanswer}
\end{myexercise}

\begin{mychallenge}
	\label{challenge:hexAddition}
	Die schriftliche Addition von Hexadezimalzahlen funktioniert analog zur schriftlichen Addition von Binärzahlen. Führen Sie die folgenden beiden Additionen durch:
	$$ 0\text{C}_{16} + 1\text{D}_{16} \qquad \text{und} \qquad \text{FF}90_{16} + 039\text{D}_{16}$$
	\begin{myanswer}
		\[
			\begin{array}{ccc}
				  & 0               & \text{C}_{16}              \\
				+ & \underset{1}{1} & \underset{}{\text{D}_{16}} \\
				\hline
				  & 2               & 9_{16}                     \\
				\hline\hline
			\end{array}
		\]

		\[
			\begin{array}{cccccc}
				  &                & \text{F}        & \text{F}        & 9              & 0_{16}                     \\
				+ & \underset{1}{} & \underset{1}{0} & \underset{1}{3} & \underset{}{9} & \underset{}{\text{D}_{16}} \\
				\hline
				  & 1              & 0               & 3               & 2              & \text{D}_{16}              \\
				\hline\hline
			\end{array}
		\]
	\end{myanswer}
\end{mychallenge}

\begin{mychallenge}\label{challenge:online}
	\begin{itemize}
		\item Wie lange benötigen Sie für 6 Umwandlungen? Testen Sie sich \href{https://oinf.ch/interactive/umrechnen-zwischen-zahlensystemen/}{hier}.
		\item Testen Sie Ihre Zahlensystem-Umwandlungs-Wissen \href{https://oinf.ch/interactive/basiswechsel-10-zu-2/}{hier}!
	\end{itemize}
\end{mychallenge}

Nachdem wir nun wissen, wie sich Zahlen in verschiedenen Stellenwertsystemen darstellen lassen, stellt sich die nächste Frage: Wie stellt man in Computern damit Zeichen, Farben, Bilder und andere Informationen dar? Diesem Thema widmet sich \cref{ch:kodierungen}.

